% 01-overview.tex — JPEG XS 概述与设计目标

\chapter{JPEG XS 概述与设计目标}\label{ch:overview}

\section{本章目标}

\begin{itemize}
  \item 理解 JPEG XS 的设计动机和应用场景
  \item 了解编码管线的整体架构
  \item 理解各模块之间的数据流转关系
\end{itemize}

\section{设计动机}

JPEG XS（ISO/IEC 21122）是面向专业视频和图像应用的轻量级压缩标准。与传统图像压缩（如 JPEG、JPEG 2000）不同，JPEG XS 的设计目标是：

\begin{itemize}
  \item \textbf{低延迟}：编码延迟在几行像素级别，不需要缓存整帧
  \item \textbf{低复杂度}：计算量适合 FPGA 和 ASIC 实现
  \item \textbf{视觉无损}：在目标压缩比（2--6$\times$）下，人眼不可区分
  \item \textbf{可逆}：支持 mathematically lossless 模式
\end{itemize}

典型应用包括：专业视频制作、机器视觉、汽车传感器、医疗影像等。

\section{编码管线}

JPEG XS 的编码管线由以下步骤组成（按执行顺序）：

\begin{enumerate}
  \item \textbf{NLT}（Non-Linear Transform）：将无符号像素映射为以零为中心的有符号系数
  \item \textbf{MCT}（Multiple Component Transform）：颜色空间变换（RCT 或 ICT）
  \item \textbf{DWT}（Discrete Wavelet Transform）：5/3 Le Gall 小波分解
  \item \textbf{Precinct 提取}：将子带系数组织为 precinct 结构
  \item \textbf{GCLI 计算}：计算每个组的全局 CLI 值
  \item \textbf{码率控制}：通过 Q/R 搜索分配预算
  \item \textbf{量化}：根据 GTLI 截断 bitplane
  \item \textbf{打包}：组装 marker 段和数据段，生成码流
\end{enumerate}

\begin{equation}
\text{RGB} \xrightarrow{\text{NLT}} \xrightarrow{\text{MCT}} \text{YUV} \xrightarrow{\text{DWT}} \text{子带} \xrightarrow{\text{Precinct}} \xrightarrow{\text{GCLI}} \xrightarrow{\text{RC}} \xrightarrow{\text{Quant}} \xrightarrow{\text{Pack}} \text{Codestream}
\end{equation}

\section{实现映射}

\begin{implmap}
\begin{tabular}{@{}L{2.0cm}L{2.0cm}L{3.0cm}L{4.5cm}@{}}
\toprule
标准概念 & 入口文件 & 关键函数/结构 & 说明 \\
\midrule
编码器总入口 & \btt{xs\_enc.c} & \btt{xs\_enc\_image()} & 单帧编码入口 \\
解码器总入口 & \btt{xs\_dec.c} & \btt{xs\_dec\_image()} & 单帧解码入口 \\
编码配置 & \btt{xs\_config.c} & \btt{xs\_config\_init()} & 参数初始化 \\
图像数据 & \btt{xs\_image.h} & \btt{xs\_image\_t} & 图像存储结构 \\
\bottomrule
\end{tabular}
\end{implmap}
